\section{Asesmen Komprehensif}

Bab ini berisi asesmen untuk mengukur pencapaian \textbf{CPMK-1} dan \textbf{CPMK-2} sesuai silabus \cite{stallings2017,munir2019}.

\subsection{Soal UTS (CPMK-2 dan Sub-CPMK 1.1--1.2 sebagian)}

UTS mengukur Sub-CPMK 2.1--2.3 serta pengenalan klasik dan simetris (1.1--1.2). Contoh soal:
\begin{enumerate}
  \item Jelaskan layanan keamanan (CIA + non-repudiation) dan urgensi kriptografi di TI (contoh HTTPS/perbankan).
  \item Hitung invers modular sederhana, misalnya $17^{-1} \pmod{26}$.
  \item Enkripsi singkat dengan Caesar/Vigen\`{e}re; jelaskan gagasan frequency analysis.
  \item Bedakan stream cipher dan block cipher; sebutkan gagasan Feistel (DES) dan SPN (AES).
  \item Bandingkan mode ECB, CBC, dan CTR secara konsep. Mengapa ECB tidak aman?
\end{enumerate}

\subsection{Soal UAS (CPMK-1 dan CPMK-2, termasuk Sub-CPMK 1.3--1.4)}

UAS mengukur seluruh jalur silabus, dengan penekanan asimetris, hash/MAC/sig, dan TLS/urgensi. Contoh soal:
\begin{enumerate}
  \item Dengan $p=5$, $q=11$, $e=3$, hitung $n$, $\varphi(n)$, $d$, dan enkripsi $m=9$ (RSA pengenalan).
  \item Jelaskan Diffie--Hellman dan ElGamal secara singkat; sebutkan masalah matematika yang mendasarinya.
  \item Bedakan hash, MAC, dan tanda tangan digital. Apa peran X.509?
  \item Jelaskan gagasan TLS/HTTPS dan mengapa validasi sertifikat penting (urgensi di TI).
  \item Bandingkan ciri kriptografi klasik vs modern serta simetris vs asimetris.
\end{enumerate}

Rubrik dan metode penilaian diuraikan pada section berikut.

